\section{Conclusion}
\label{sec:conclusion}

This work has formulated horizon occlusion as an inverse problem in
geometrical optics.  Starting from Fermat's principle for a horizontally
stratified refractive index $n(h)$, we derived the ray equation in both the
paraxial and exact regimes and asked which vertical profile would reproduce
the limiting visibility law associated with a curved surface.

In the paraxial approximation, constant upward ray curvature requires an
exponential refractive index,
\begin{equation}
    n(h)=n_0 e^{h/R},
\end{equation}
when the curvature scale is chosen to match a sphere of radius $R$.  The
resulting grazing rays recover the familiar leading horizon law
\begin{equation}
    d(h)\simeq\sqrt{2Rh},
\end{equation}
and the maximum observer--target separation is the sum of the two one-sided
horizon distances.

The exponential model is also exactly solvable.  Its ray family is
\begin{equation}
    h(x)
    =
    h_m
    -R\ln\!\left[
        \cos\!\left(\frac{x-x_m}{R}\right)
    \right],
\end{equation}
and its exact one-sided horizon distance is
\begin{equation}
    d_{\mathrm e}(h)
    =
    R\arccos\!\left(e^{-h/R}\right).
\end{equation}
A conformal coordinate transformation maps the corresponding optical metric
to the homogeneous Euclidean exterior of a circle.  This establishes an exact
optical equivalence between the flat exponential medium and circular
occlusion, with the important qualification that physical height is mapped
nonlinearly according to $r=Re^{h/R}$.

Comparison with the exact spherical arc-distance law,
\begin{equation}
    d_{\mathrm s}(h)
    =
    R\arccos\!\left(\frac{R}{R+h}\right),
\end{equation}
shows that the two models agree extremely well at terrestrial heights but are
not identical at equal numerical height beyond leading order.  Exact equality
for a prescribed horizon law can instead be enforced through the general
inverse reconstruction
\begin{equation}
    n(h)
    =
    n_0\sqrt{1+\frac{1}{[d'(h)]^2}},
\end{equation}
within the assumptions of a smooth monotonic stratified medium and a single
grazing branch.

The decisive physical result is the sign of the required gradient.  Both the
exponential construction and the profile reconstructed from the exact
spherical law require
\begin{equation}
    \frac{dn}{dh}>0,
\end{equation}
so that rays bend upward and possess lower turning points.  Ordinary
atmospheric stratification has the opposite sign, $dn/dh<0$, apart from local
and transient anomalous layers.  Standard atmospheric refraction therefore
cannot serve as a global mechanism that reproduces the complete spherical
horizon law above a flat boundary.

The analysis nevertheless demonstrates a broader point.  Bottom-first
disappearance, recovery with observer height, and a finite horizon are not
logically unique signatures of surface curvature when arbitrary refractive
media are permitted.  They can be reproduced by a specially engineered
optical geometry.  What distinguishes competing explanations is not the mere
existence of bent rays, but the quantitative refractive profile, its physical
stability, and the additional imaging effects it predicts.  In that sense,
the inverse problem converts a qualitative resemblance into a testable
condition on the medium.
